% !TeX program = xelatex
% !TEX encoding = UTF-8
\documentclass[11pt,a4paper]{article}

\usepackage{zhpl}

% ---------- Обозначения ----------
\newcommand{\Rfield}{\mathcal{R}}
\newcommand{\kB}{k_{\text{Б}}}
\newcommand{\Burm}{\text{Б}}
\newcommand{\Lnar}{\Lambda_{\text{д}}}

\title{\textbf{К квантовополевому обоснованию рояля в кустах:\\ вакуумная нуклеация, казимировское условие на кустарник и драматургическая поправка}}
\author{\zhplauthors}
\zhplsetup{4}{10.9999/жпл.2026.0004}

\begin{document}

\maketitle

\begin{center}
\begin{tikzpicture}
\node[draw=\zhplbadgecolor{red!70!black}{zhplBadgeRed},very thick,rounded corners,rotate=-6,fill=\zhpldark{red!5},inner sep=5pt] (badge1) {\color{\zhplbadgecolor{red!70!black}{zhplBadgeRed}}\bfseries ПРОРЫВНОЕ ИССЛЕДОВАНИЕ};
\node[draw=\zhplbadgecolor{blue!60!black}{zhplBadgeBlue},very thick,rounded corners,rotate=4,fill=\zhpldark{blue!5},inner sep=5pt,right=1.0cm of badge1] (badge2) {\color{\zhplbadgecolor{blue!60!black}{zhplBadgeBlue}}\bfseries IMPACT FACTOR: $\infty$};
\end{tikzpicture}
\end{center}

\zhplciteas{Александров Р.\,Г., Антропикович К.\,С. (2026). К квантовополевому обоснованию рояля в кустах.
  \textit{Журнал Прикладной Лженауки}, \textbf{1}(4), 1--9. DOI: 10.9999/жпл.2026.0004}

\zhplgraphicabstract{%
  \begin{tikzpicture}[node distance=0.75cm, every node/.style={font=\footnotesize}]
    \node[draw,rounded corners,fill=\zhpldark{blue!5},minimum width=2.4cm,minimum height=1.05cm,align=center] (n1) {Вакуум\\(рояль виртуален)};
    \node[draw,rounded corners,fill=\zhpldark{blue!5},minimum width=2.5cm,minimum height=1.05cm,align=center,right=of n1] (n2) {Куст\\(резонатор, $d=nL/2$)};
    \node[draw,rounded corners,fill=\zhpldark{blue!5},minimum width=2.5cm,minimum height=1.05cm,align=center,right=of n2] (n3) {Сюжет\\(кривизна $\Xi$)};
    \node[draw,rounded corners,fill=\zhpldark{green!8},minimum width=2.4cm,minimum height=1.05cm,align=center,right=of n3] (n4) {Рояль\\(строго вовремя)};
    \draw[-{Latex[length=2.5mm]},very thick] (n1) -- (n2);
    \draw[-{Latex[length=2.5mm]},very thick] (n2) -- (n3);
    \draw[-{Latex[length=2.5mm]},very thick] (n3) -- (n4);
  \end{tikzpicture}}{%
  От вездесущей виртуальной флуктуации к наблюдаемому изделию: куст задаёт граничное условие, сюжет~--- показатель экспоненты, ценой всякой возможности проверить это намеренно.}

\begin{zhplhighlights}
\begin{itemize}
\setlength\itemsep{0pt}
\setlength\parskip{0pt}
\item Впервые показано, что рояль в кустах~--- не литературный приём, а распад ложного безрояльного вакуума.
\item Доказано (Теорема~1), что рояль присутствует в любом кусте в любой момент времени; наблюдению препятствует исключительно время жизни $2{,}4\cdot10^{-54}$~с.
\item Установлено казимировское условие на кустарник: нуклеация резонансно усилена при межкустовом расстоянии, кратном $1{,}37$~м.
\item Сформулировано правило отбора по тембральному заряду, из которого следует парное рождение «рояль\,+\,пианино» (Теорема~3) и свободное рождение баянов.
\item Доказано (Теорема~4), что вероятность обнаружения рояля целенаправленным поиском строго меньше вероятности его обнаружения при отсутствии поиска.
\end{itemize}
\end{zhplhighlights}

\begin{zhplreviews}
\textit{«До прочтения я двадцать лет считал рояль в кустах недостатком сценария. Теперь я понимаю, что это единственный эпизод, физику которого автор просчитал корректно.»} -- Рецензент 1\\
\textit{«Замечаний нет.»} -- Рецензент 2\\
\textit{«Согласно Теореме~4 настоящей статьи, целенаправленное рецензирование обнуляет нарративную кривизну текста, вследствие чего рояль в нём заведомо не появится. Отзыв составлен без такового и потому объективным считаться не может.»} -- Рецензент 3
\end{zhplreviews}

\begin{zhplsignificance}
Наивная (нефизическая) литературоведческая традиция полтора столетия относила рояль в кустах к числу авторских произволов~--- явлений, по определению не поддающихся количественному описанию и потому исключённых из сферы науки. Настоящая работа показывает, что это не так: появление рояля в кустах есть распад метастабильного состояния вакуума, скорость которого вычисляется из первых принципов, а его наблюдаемая привязка к сюжетной необходимости~--- не свидетельство произвола, а прямое следствие вида показателя экспоненты. Результат применим к любому кустарнику на планете, а также, с точностью до замены изделия, ко всей художественной литературе, накопленной человечеством.
\end{zhplsignificance}

\begin{abstract}
\noindent Здесь мы впервые показываем, что рояль в кустах~--- не драматургическая условность, а строгий квантовополевой эффект: распад ложного безрояльного вакуума с образованием пузыря новой фазы. Бытовое объяснение («автор не придумал ничего лучше») игнорирует базовый факт: событие не запрещено ни одним законом сохранения, а следовательно, обязано происходить. В настоящей работе мы вводим рояльное поле $\Rfield$ и его двухъямный потенциал, вычисляем скорость нуклеации в подходе ложного вакуума, показываем, что виртуальный рояль присутствует в каждом кусте постоянно и что задача состоит не в его создании, а в стабилизации, устанавливаем роль кустарника как казимировского резонатора с резонансным условием $d = nL_{\text{рояль}}/2$, формулируем правило отбора по тембральному заряду и вводим драматургическую поправку $\Lnar\Xi$ к показателю экспоненты, связывающую скорость нуклеации с нарративной кривизной сюжета. Приводятся результаты полевых измерений на выборке из 1174 кустов. Обсуждаются ограничения метода.

\medskip
\noindent\textbf{Ключевые слова:} ложный вакуум, нуклеация, эффект Казимира, правила отбора, тембральный заряд, нарративная кривизна, лженаука
\end{abstract}

\section{Введение}

Феномен рояля в кустах известен не менее полутора столетий и до сих пор не имеет микроскопического описания. Эмпирически установлено следующее: рояль обнаруживается в кустах ровно тогда, когда он там нужен; никогда не обнаруживается при целенаправленном поиске; и практически никогда не оказывается расстроенным. Последнее обстоятельство особенно примечательно, поскольку исключает гипотезу о том, что рояль был занесён в кусты заранее: инструмент, простоявший в кустарнике хотя бы одну ночь, теряет строй необратимо.

Литературоведческая традиция объясняет всё это авторским произволом. Такое объяснение не содержит ни одной дифференцируемой величины и потому научной ценности не имеет. Мы исходим из противоположной посылки: если событие не запрещено законами сохранения, оно происходит, и вопрос состоит лишь в том, с какой скоростью и при каких граничных условиях. Ниже показано, что задача о рояле в кустах полностью укладывается в стандартную теорию распада метастабильного вакуума, а вся её кажущаяся «литературность» сосредоточена ровно в одном слагаемом показателя экспоненты.

\section{Рояльное поле и ложный безрояльный вакуум}

\begin{definition}[Рояльное поле]
Рояльным полем называется вещественное скалярное поле $\Rfield(x)$ с действием
\[
S[\Rfield] = \int d^4x \left[ \tfrac{1}{2}\,\partial_\mu \Rfield\,\partial^\mu \Rfield - V(\Rfield) \right],
\]
потенциал которого имеет два локальных минимума: $\Rfield = \Rfield_\emptyset$ (кусты пусты) и $\Rfield = \Rfield_0$ (кусты содержат рояль), разделённых барьером конечной высоты.
\end{definition}

Наблюдаемая нами повседневность соответствует состоянию $\Rfield_\emptyset$, которое, вопреки бытовому впечатлению, не является основным. Оно метастабильно: рояльная фаза энергетически выгоднее, поскольку разрешает больше сюжетов. Скорость распада на единицу объёма даётся стандартным выражением теории ложного вакуума
\begin{equation}
\frac{\Gamma}{V} = A\,\exp\!\left(-\frac{S_{\mathrm{E}}}{\hbar}\right),
\label{eq:gamma}
\end{equation}
где $S_{\mathrm{E}}$~--- евклидово действие критического пузыря, а предэкспоненциальный множитель $A$ имеет размерность $\text{рояль}\cdot\text{м}^{-3}\text{с}^{-1}$ и далее полагается равным единице в силу отсутствия независимых измерений.

Оценка показателя экспоненты через энтропию изделия даёт величину, ради которой стоит завести отдельную строку. Для рояля массой $M = 480$~кг, состоящего из палисандра, чугуна и надежд, число степеней свободы составляет $N\sim10^{28}$, откуда
\begin{equation}
P \sim \exp\!\left(-\frac{S_{\text{рояль}}}{\kB}\right) \sim 10^{-4{,}34\cdot10^{27}}.
\label{eq:naive}
\end{equation}
Характерное время ожидания превышает возраст Вселенной настолько, что разница между «превышает» и «равно» лежит за пределами точности любых мыслимых измерений. Обычно на этом месте исследование прекращают.

\begin{figure}[htbp]
\centering
\begin{tikzpicture}
\begin{axis}[
  width=11cm, height=6cm,
  xlabel={$\Rfield$}, ylabel={$V(\Rfield)$},
  domain=-1.55:1.55, samples=201,
  ymin=-0.35, ymax=0.75,
  axis lines=left, grid=major,
  legend pos=north west, legend cell align=left,
]
\addplot[very thick,blue] {0.25*(x^2-1)^2 - 0.15*x};
\addlegendentry{$V(\Rfield)=\tfrac14(\Rfield^2-1)^2-\varepsilon\Rfield$, $\varepsilon=0{,}15$}
\addplot[only marks,mark=*,red] coordinates {(-1.0742,0.1780)};
\addplot[only marks,mark=*,green!50!black] coordinates {(1.0742,-0.1443)};
\addplot[only marks,mark=x,mark size=4pt,black] coordinates {(-0.1483,0.2670)};
\node[anchor=west,font=\footnotesize] at (axis cs:-1.50,0.45) {ложный вакуум: кусты пусты};
\draw[-{Latex[length=1.6mm]},thin] (axis cs:-1.22,0.40) -- (axis cs:-1.08,0.21);
\node[anchor=east,font=\footnotesize] at (axis cs:1.50,-0.30) {истинный вакуум: рояль};
\draw[-{Latex[length=1.6mm]},thin] (axis cs:1.18,-0.27) -- (axis cs:1.08,-0.17);
\node[anchor=south,font=\footnotesize] at (axis cs:-0.1483,0.2900) {барьер $S_{\mathrm{E}}$};
\end{axis}
\end{tikzpicture}
\caption{Потенциал рояльного поля. Точки минимумов и максимума получены решением $V'(\Rfield)=\Rfield^3-\Rfield-\varepsilon=0$ при $\varepsilon=0{,}15$: $\Rfield=-1{,}0742$; $-0{,}1483$; $1{,}0742$. Разность глубин $2\varepsilon\Rfield_0 = 0{,}322$ соответствует выигрышу в числе доступных сюжетов.}
\label{fig:potential}
\end{figure}

\section{Вездесущность рояля}

Даже состоявшаяся флуктуация ограничена соотношением неопределённостей $\Delta E\,\Delta t\sim\hbar$.

\begin{theorem}[О вездесущности рояля]
В любом кусте в любой момент времени рояль присутствует. Время его существования равно
\[
\tau = \frac{\hbar}{Mc^2},
\]
что для концертного рояля массой $480$~кг составляет $2{,}4\cdot10^{-54}$~с.
\end{theorem}

Этот результат принято понимать неправильно. Он означает не то, что роялей в кустах не бывает, а то, что они там \emph{всегда}: виртуальная рояльная компонента вакуума ненулевая в любом объёме, включая рассматриваемый. Проблема исключительно наблюдательная. За $2{,}4\cdot10^{-54}$~с исполнитель не успевает не только сыграть, но и сесть; более того, за это время свет проходит расстояние порядка $10^{-45}$~м, то есть рояль не успевает стать видимым даже для собственных струн.

Таким образом, задача переформулируется: требуется не \emph{создать} рояль, а \emph{стабилизировать} уже имеющийся. Заметим, что $\tau\propto M^{-1}$ (рис.~\ref{fig:tau}), откуда немедленно следует известный полевой факт: губная гармоника из вакуума достаётся втрое легче, чем баян, и на четыре порядка легче рояля.

\begin{figure}[htbp]
\centering
\begin{tikzpicture}
\begin{axis}[
  width=11cm, height=6.2cm,
  xmode=log, ymode=log,
  xlabel={масса изделия $M$, кг}, ylabel={$\tau=\hbar/Mc^2$, с},
  domain=0.03:800, samples=60,
  grid=major,
  legend pos=north east, legend cell align=left,
]
\addplot[very thick,blue] {1.1734e-51/x};
\addlegendentry{$\tau = \hbar/(Mc^2)$}
\addplot[only marks,mark=*,red] coordinates {
  (0.05,2.347e-50) (2,5.867e-52) (12,9.778e-53) (250,4.693e-54) (480,2.445e-54)};
\addlegendentry{изделия}
\node[anchor=west,font=\footnotesize] at (axis cs:0.05,2.347e-50) {\ гармоника};
\node[anchor=west,font=\footnotesize] at (axis cs:2,5.867e-52) {\ гитара};
\node[anchor=west,font=\footnotesize] at (axis cs:12,9.778e-53) {\ баян};
\node[anchor=east,font=\footnotesize] at (axis cs:250,4.693e-54) {пианино\ };
\node[anchor=north east,font=\footnotesize] at (axis cs:430,2.0e-54) {рояль\ };
\end{axis}
\end{tikzpicture}
\caption{Время жизни виртуального музыкального изделия в зависимости от его массы; $\hbar/c^2 = 1{,}1734\cdot10^{-51}$~кг$\cdot$с. Точки вычислены по той же формуле, а не расставлены по месту: гармоника $0{,}05$~кг~--- $2{,}35\cdot10^{-50}$~с, рояль $480$~кг~--- $2{,}45\cdot10^{-54}$~с.}
\label{fig:tau}
\end{figure}

\section{Кусты как граничное условие}

До сих пор кусты рассматривались как декорация. Это ошибка. Куст есть периодическая диэлектрическая структура, налагающая на вакуумные моды граничные условия, то есть казимировский резонатор с листьями.

\begin{definition}[Кустистость]
Кустистостью $B$ называется плотность ветвления кустарника, определяемая как число отражающих элементов на единицу объёма. Куст с $B \to 0$ называется газоном.
\end{definition}

Два куста, разделённые расстоянием $d$, вырезают из вакуумного спектра моды $\lambda_n = 2d/n$. Если собственная мода изделия попадает в разрешённую полосу, соответствующая флуктуация усиливается; коэффициент усиления имеет стандартный вид резонатора Фабри--Перо
\begin{equation}
G(d) = \left[1 + F\sin^2\!\left(\frac{2\pi d}{L_{\text{рояль}}}\right)\right]^{-1},
\qquad L_{\text{рояль}} = 2{,}74\ \text{м},
\label{eq:airy}
\end{equation}
и обращается в максимум при
\begin{equation}
d = \frac{n}{2}\,L_{\text{рояль}} = 1{,}37n\ \text{м}, \qquad n = 1,2,3,\dots
\label{eq:resonance}
\end{equation}

Это первое количественное предсказание теории: рояли возникают преимущественно в кустах, высаженных с шагом, кратным $1{,}37$~м. Известная садоводческая рекомендация сажать кустарник «примерно через полтора метра» получает тем самым глубокое физическое обоснование, а вся ландшафтная архитектура задним числом переводится в разряд прикладной квантовой теории поля.

\begin{figure}[htbp]
\centering
\begin{tikzpicture}
\begin{axis}[
  width=11cm, height=5.6cm,
  xlabel={межкустовое расстояние $d$, м}, ylabel={$G(d)$},
  domain=0.05:4.5, samples=600,
  ymin=0, ymax=1.15, xmin=0, xmax=4.5,
  grid=major, legend pos=north east, legend cell align=left,
  xtick={0,1.37,2.74,4.11}, xticklabels={$0$,$1{,}37$,$2{,}74$,$4{,}11$},
]
\addplot[very thick,blue] {1/(1+50*(sin(deg(2*pi*x/2.74)))^2)};
\addlegendentry{$F=50$}
\addplot[dashed,red,domain=0:4.5,samples=2] {1};
\end{axis}
\end{tikzpicture}
\caption{Коэффициент усиления рояльной моды по формуле~\eqref{eq:airy}: настоящая функция Эйри резонатора, а не нарисованные пики. Максимумы приходятся ровно на $d = 1{,}37n$~м; при произвольной посадке кустарника усиление составляет в среднем $1/\sqrt{1+F} \approx 0{,}14$ от резонансного.}
\label{fig:airy}
\end{figure}

Кустарник решает и вторую задачу~--- удержания пузыря новой фазы.

\begin{theorem}[О редком кустарнике]
Пузырь рояльной фазы радиуса $R$ расширяется только при $R > R_c = 2\sigma/\varepsilon$, где $\sigma$~--- поверхностное натяжение стенки пузыря, $\varepsilon$~--- разность плотностей энергии фаз. Поскольку куст вносит в баланс дополнительный вклад $-\,\beta B$, снижающий эффективное $\sigma$, при $B \to 0$ имеем $R_c \to \infty$, и рояль в газоне не удерживается ни при каком радиусе.
\end{theorem}

Полевые наблюдения (раздел~7) полностью подтверждают этот вывод: за всё время исследования в газоне не было обнаружено ни одного рояля.

\section{Правила отбора}

Возникает естественный вопрос: почему рояль, а не пианино, и почему не, скажем, ксилофон.

\begin{definition}[Тембральный заряд]
Тембральным зарядом $T$ называется чётность инструмента относительно оси струнодержателя:
\[
T(\text{рояль}) = +1, \qquad T(\text{пианино}) = -1, \qquad T(\text{баян}) = 0.
\]
\end{definition}

\begin{theorem}[Правило отбора]
Вакуум обладает нулевым тембральным зарядом. Следовательно, рождение изделия с $T=\pm1$ поодиночке запрещено и возможно только парами с суммарным $T=0$; рождение изделий с $T=0$ не ограничено.
\end{theorem}

Отсюда следует второе предсказание: \emph{рояль в кустах всегда сопровождается пианино в соседних кустах}. Проверок этого утверждения в литературе не обнаружено, что объясняется не сложностью эксперимента, а тем, что в соседние кусты никто не смотрел. Отдельно отметим, что баян, обладая нулевым зарядом, рождается свободно и в неограниченных количествах~--- вывод, согласующийся как с теорией, так и с бытовым опытом.

\section{Драматургическая поправка}

Ни казимировский механизм, ни правила отбора не объясняют главного: корреляции между появлением рояля и сюжетной необходимостью.

\begin{definition}[Нарративная кривизна]
Нарративной кривизной $\Xi$ называется свёртка тензора личности персонажа~\cite{tenzor} с метрикой сюжета в точке события. Величина $\Xi$ безразмерна и нормирована так, что $\Xi = 1$ отвечает эпизоду, без которого повествование не может быть продолжено.
\end{definition}

Модифицированная формула нуклеации принимает вид
\begin{equation}
\boxed{\;\frac{\Gamma}{V} = A\,\exp\!\left(-\frac{S_{\mathrm{E}}}{\hbar} + \Lnar\,\Xi\right)\;}
\label{eq:master}
\end{equation}
где $\Lnar \approx 10^{28}$~--- драматургическая постоянная, подобранная так, чтобы при $\Xi = 1$ показатель экспоненты обращался в нуль. Столь неизящный подбор параметра является, по нашему мнению, не недостатком теории, а её главным достижением: он объясняет, почему рояль появляется не «часто» и не «редко», а строго вовремя (рис.~\ref{fig:master}).

\begin{figure}[htbp]
\centering
\begin{tikzpicture}
\begin{axis}[
  width=11cm, height=5.8cm,
  xlabel={нарративная кривизна $\Xi$}, ylabel={$\log_{10} P$, единиц $10^{27}$},
  domain=0:2, samples=2,
  ymin=-5, ymax=5, xmin=0, xmax=2,
  grid=major, legend pos=north west, legend cell align=left,
]
\addplot[very thick,blue] {4.3429*(x-1)};
\addlegendentry{$\log_{10}P = \Lnar(\Xi-1)/\ln 10$}
\addplot[dashed,red,domain=0:2,samples=2] {0};
\addlegendentry{$P=1$}
\addplot[only marks,mark=*,black] coordinates {(1,0)};
\node[anchor=north west,font=\footnotesize] at (axis cs:1.05,-0.2) {порог: рояль появляется};
\node[anchor=south east,font=\footnotesize] at (axis cs:1.95,-4.8) {экспедиция за роялем: $\Xi=0$};
\end{axis}
\end{tikzpicture}
\caption{Вероятность нуклеации по формуле~\eqref{eq:master} при $S_{\mathrm{E}}/\hbar = \Lnar = 10^{28}$. Прямая построена по точной формуле $\log_{10}P = 10^{28}(\Xi-1)/\ln 10 = 4{,}3429\cdot10^{27}(\Xi-1)$. При $\Xi = 0$ имеем $P = 10^{-4{,}34\cdot10^{27}}$, что совпадает с наивной оценкой~\eqref{eq:naive}; переход через единицу происходит на интервале $\Delta\Xi \sim 10^{-28}$, то есть мгновенно.}
\label{fig:master}
\end{figure}

Из~\eqref{eq:master} следует \emph{принцип наименьшего действия сценариста}: рояль материализуется в той точке пространства-времени, где вариация нарративного действия обращается в нуль, то есть там, где его появление меньше всего портит историю.

\begin{theorem}[О ненаблюдаемости при поиске]
Пусть наблюдатель ведёт целенаправленный поиск рояля в кустах. Тогда его обнаружение не является сюжетно необходимым (эпизод продолжится и без него), откуда $\Xi \to 0$ и $P \to 10^{-4{,}34\cdot10^{27}}$. Следовательно, вероятность обнаружить рояль при поиске строго меньше вероятности обнаружить его при отсутствии поиска.
\end{theorem}

Теорема~4 объясняет всю имеющуюся отрицательную статистику, включая нашу собственную (раздел~7), и делает теорию нефальсифицируемой~--- обстоятельство, которое мы относим к её достоинствам.

Наконец, поскольку нарративная кривизна пропорциональна градиенту бурмалды~\cite{burmalda},
\begin{equation}
\Xi \propto \frac{\partial\,\Burm}{\partial t},
\end{equation}
скорость нуклеации максимальна в момент наиболее резкого роста бурмалды, то есть примерно через три часа после начала мероприятия. Этим объясняется фундаментальный факт: рояль в кустах обнаруживается почти исключительно на дачах и почти исключительно ночью.

\section{Экспериментальная часть}

В период с мая по август обследовано 1174 куста (Флеволанд, Нидерланды; преимущественно бирючина и боярышник). Методика: визуальный осмотр, простукивание, прослушивание на предмет спонтанного \emph{до} первой октавы. Контрольная группа~--- 40~м$^2$ газона. Результаты сведены в таблицу~\ref{tab:field}.

\begin{table}[htbp]
\centering
\begin{tabular}{lcc}
\toprule
Объект & Обнаружено & Предсказано теорией \\
\midrule
Рояль концертный        & 0  & $\ll 1$ \\
Рояль кабинетный        & 0  & $\ll 1$ \\
Пианино                 & 0  & $= N_{\text{рояль}}$ \\
Рок-группа репетирующая & 1  & не рассматривалась \\
Баян                    & 1  & $\gg 1$ \\
Кот                     & 17 & вне компетенции \\
\bottomrule
\end{tabular}
\caption{Результаты полевых измерений. Выборка: 1174 куста, контроль~--- газон (находок нет).}
\label{tab:field}
\end{table}

Согласие с теорией следует признать превосходным. Нулевое число роялей в точности воспроизводит предсказание~\eqref{eq:naive} при $\Xi = 0$, гарантированном самим фактом проведения поиска (Теорема~4); обнаруженный баян подтверждает правило отбора по тембральному заряду (Теорема~3); равенство числа пианино числу роялей выполнено с абсолютной точностью. Репетирующая рок-группа теорией не предсказывалась, но и не запрещается: суммарный тембральный заряд коллектива равен нулю, а собственная бурмалда состава достаточна для самостабилизации пузыря без участия куста. Именно этим объясняется то, что данное событие было зарегистрировано с расстояния 300~м и оказалось единственным за всё исследование, не потребовавшим простукивания. Семнадцать котов отнесены к фоновым событиям.

\section{Обсуждение и ограничения}

Метод обладает тремя ограничениями, которые мы считаем необходимым назвать прямо.

\emph{Первое.} Драматургическая постоянная $\Lnar$ подобрана из единственного требования~--- чтобы ответ получился правильным. Независимого измерения $\Lnar$ не существует и, ввиду Теоремы~4, существовать не может: любая процедура измерения есть целенаправленный поиск и обнуляет измеряемую величину.

\emph{Второе.} При $\Xi > 1$ формула~\eqref{eq:master} даёт $P > 1$. Стандартная интерпретация~--- нарушение унитарности~--- в нашем случае неприменима, поскольку означала бы, что рояль появляется чаще, чем всегда. Мы предлагаем считать, что при $\Xi > 1$ рояль появляется независимо от наличия кустов, вакуума и физики, что согласуется с наблюдениями.

\emph{Третье.} Теория не даёт ни малейшего указания на то, кто настраивает инструмент. Между тем именно настроенность рояля составляет главную эмпирическую загадку явления и единственную его деталь, которую невозможно объяснить флуктуацией: расстроенный рояль имеет большую энтропию, а значит, и заведомо большую вероятность нуклеации, чем настроенный. Наблюдается обратное. Поскольку теория предсказывает преимущественное появление именно расстроенных роялей, а наблюдаются исключительно настроенные, настоящая работа опровергнута собственным основным предсказанием. Это можно рассматривать как единственный строгий результат работы.

\section{Заключение}

Мы показали, что спонтанное возникновение рояля в кустах не противоречит ни одному известному закону физики, а следовательно, обязательно; что вероятность события ничтожна, но конечна, а время жизни изделия ничтожно, но ненулевое; что кустарник выполняет функцию казимировского резонатора с резонансным условием $d = 1{,}37n$~м; и что драматургическая поправка приводит показатель экспоненты к значениям порядка единицы ровно тогда, когда это требуется по сюжету. Открытым остаётся вопрос о настройке инструмента, для решения которого требуется постановка отдельного эксперимента с существенно большей выборкой кустов. Мы считаем это достаточным основанием для дальнейшего финансирования настоящего направления исследований.

\section*{Вклад авторов}
Р.\,Г.~Александров: концептуализация, методология, полевые измерения (осмотр и простукивание 1174 кустов), написание -- черновик.
К.\,С.~Антропикович: программное обеспечение, визуализация, формальный анализ, написание -- рецензирование и редактирование, техническая поддержка бесконечного терпения.

\section*{Финансирование}
Исследование выполнено при поддержке гранта Института Прикладной Метафизики
№ 2026-$\Omega$-$\infty$ <<Инвентаризация кустарников общего пользования и сопутствующих метастабильных состояний>>.

\section*{Доступность данных}
Данные доступны по обоснованному запросу -- а также по необоснованному, с той же вероятностью успеха. Кусты доступны по адресу проведения измерений, за исключением семнадцати, отозвавших согласие на участие.

\begin{thebibliography}{9}
\bibitem{tenzor} Александров, Р.\,Г., Антропикович, К.\,С. \textit{К обоснованию тензорного анализа личности: инвариантная психометрика в искривлённых пространствах характера.} Журнал Прикладной Лженауки, \textbf{1}(1), 1--7 (2026).
\bibitem{burmalda} Александров, Р.\,Г., Антропикович, К.\,С. \textit{К термодинамическому обоснованию бурмалды: экстенсивная величина, монотонно неубывающая в замкнутой компании.} Журнал Прикладной Лженауки, \textbf{1}(3), 1--9 (2026).
\bibitem{kustovoy} Кустовой, К.\,К. \textit{Предельная кустистость живой изгороди и условия удержания пузыря новой фазы.} Труды Института Прикладной Метафизики, б.г.
\bibitem{tembrovskiy} Тембровский, Т.\,З. \textit{Правила отбора для струнных изделий, рождающихся из вакуума.} Неопубликованная рукопись.
\bibitem{nukleatov} Нуклеатов, Н.\,Н. \textit{Критический радиус пузыря в условиях садового участка: опыт шести сезонов.} Труды Института Прикладной Метафизики, б.г.
\end{thebibliography}

\noindent\textbf{Конфликт интересов.} Не заявлен, хотя, согласно Теореме~4, его отсутствие строго доказать невозможно: любая проверка наличия конфликта является целенаправленным поиском и потому заведомо ничего не обнаруживает.

\end{document}
